\chapter{Libreria \texttt{fault}}
%--------1---------2---------3---------4---------5---------6---------7---------8
La libreria contiene delle eccezioni specifiche e la definizione dei codici di
errore.

%--------1---------2---------3---------4---------5---------6---------7---------8
\begin{elisting}[!htb]
\begin{javacode}
public class Fixed {
    public static final int RC_OK       =  0;
    public static final int RC_WARN     =  4;
    public static final int RC_ERR_IO   =  8;
    public static final int RC_ERR_SQL  = 12;
    public static final int RC_ERR_STEP = 16;
    public static final int RC_ERR_JOB  = 20;
    public static final IntBinaryOperator MAX_CC = Math::max;
}
\end{javacode}
\caption{Codici di errore, valori di default}
\label{lst:rc-def}
\end{elisting}
%--------1---------2---------3---------4---------5---------6---------7---------8

%--------1---------2---------3---------4---------5---------6---------7---------8
I codici di errore sono definiti mediante una service provider interface (SPI)
in questo modo fornendo una definizione custom dell'\,interfaccia è possibile
ridefinire i valori dei codici di errore.

%--------1---------2---------3---------4---------5---------6---------7---------8
\begin{elisting}[!htb]
\begin{javacode}
public interface ValueProvider {
    ValueFactory getInstance();
}
\end{javacode}
\caption{Interfaccia per fornire il generatore del valore degli errori}
\label{lst:val-prov}
\end{elisting}
%--------1---------2---------3---------4---------5---------6---------7---------8


%--------1---------2---------3---------4---------5---------6---------7---------8
\begin{elisting}[!htb]
\begin{javacode}
public interface ValueFactory {
    int rcOk();
    int rcWarning();
    int rcErrorIO();
    int rcErrorSQL();
    int rcErrorStep();
    int rcErrorJob();
    int rcMax(int a, int b);
}
\end{javacode}
\caption{Interfaccia per fornire il valore dei codici di errore}
\label{lst:val-fact}
\end{elisting}
%--------1---------2---------3---------4---------5---------6---------7---------8
